\chapter{Resultados experimentales}\label{cap:resultados}

En este capítulo se presentan y discuten los resultados experimentales obtenidos.

\section{Experimento 1: fusión vs. modelos individuales}\label{sec:resultados-e1}
% - Tabla principal de modelos base y fusiones
% - Qué fusiones superan a qué modelos, con significancia
% - Lectura por familia (léxico vs. semántico)

\section{Experimento 2: comparación de algoritmos de fusión}\label{sec:resultados-e2}
% - Tabla comparativa de los seis algoritmos
% - Ranking de algoritmos y diferencias significativas
% - Sensibilidad a parámetros (k de RRF) si aplica

\section{Experimento 3: fusión vs. reordenamiento neuronal}\label{sec:resultados-e3}
% - Tabla: mejor fusión vs. rerankers
% - Costo/beneficio: relevancia ganada vs. costo
% - Significancia

\section{Experimento 4: modelos abiertos vs. propietarios}\label{sec:resultados-e4}
% - Tabla vs. líneas base propietarias
% - Posición relativa alcanzada

\section{Análisis de complementariedad léxico-semántica}\label{sec:resultados-complementariedad}
% - Medida de complementariedad: RBO (Rank-Biased Overlap; Webber, Moffat y Zobel, 2010) — misma familia de pesos que RBC
% - Por qué no Kendall τ: exige listas conjuntas y no pondera por profundidad; RBO admite listas no conjuntas (lo habitual entre familias) y pondera el tope
% - Consultas donde gana cada familia, con ejemplos
% - Relación entre complementariedad y ganancia de fusión (responde al objetivo específico 6)
